% ============================================================================
% CV Youssef EL KATTOUFI - candidature Bell 2027, novembre 2026
% LuaLaTeX recommande (Overleaf > Menu > Compilateur > LuaLaTeX)
% Mise a jour personnalisee a partir du CV source et de l'offre Bell 432858.
% Les objectifs sont nommes comme tels ; aucune performance non mesuree n'est presentee comme acquise.
% Identite : bleu / rouge inspires du logo Sorbonne Universite ; gris du
%           langage graphique Polytech. Aucune reproduction de logo officiel.
% Illustrations entierement vectorielles TikZ, dans la marge droite.
% ============================================================================
\documentclass[9pt,a4paper]{extarticle}
\usepackage[a4paper,top=0.95cm,bottom=0.82cm,left=1.24cm,right=1.62cm]{geometry}
\usepackage{fontspec}
\setmainfont{Carlito}[BoldFont={Carlito Bold}, ItalicFont={Carlito Italic},Ligatures=NoCommon]
\usepackage[french]{babel}
\usepackage{microtype}
\usepackage{xcolor}
\usepackage{tikz}
\usetikzlibrary{calc,shapes.geometric}
\usepackage{eso-pic}
\usepackage{tabularx}
\usepackage{array}
\usepackage{enumitem}
\usepackage{titlesec}
\usepackage[hidelinks]{hyperref}
\hypersetup{pdftitle={CV - Youssef EL KATTOUFI - Electronique et Informatique},pdfauthor={Youssef EL KATTOUFI},pdfsubject={Candidature au programme de stages 2027 en innovation technologique chez Bell},colorlinks=true,urlcolor=bleu}
\definecolor{bleu}{HTML}{263068}
\definecolor{rouge}{HTML}{E6332A}
\definecolor{gris}{HTML}{9EA6B6}
\definecolor{anthracite}{HTML}{222936}
\definecolor{secondaire}{HTML}{475269}
\definecolor{fond}{HTML}{F0F2F7}
\color{anthracite}
\pagestyle{empty}
\setlength{\parindent}{0pt}
\setlength{\parskip}{0pt}
\setlength{\emergencystretch}{1em}
\renewcommand{\arraystretch}{0.98}
\newcommand{\stage}{STAGE BELL 2027 -- DISPONIBLE DÈS AVRIL 2027}

% --- Iconographie TikZ originale, placee hors du bloc de texte --------------
\newcommand{\decorationsCV}{%
\begin{tikzpicture}[remember picture, overlay]
  % Fine ligne de connexion dans la marge droite
  \draw[bleu, opacity=.11,line width=.48pt] ([xshift=-8.0mm,yshift=-50mm]current page.north east) -- ([xshift=-8.0mm,yshift=24mm]current page.south east);
  % Puce + pistes de microcontroleur
  \begin{scope}[shift={([xshift=-8.0mm,yshift=-65mm]current page.north east)},scale=.70,opacity=.29,draw=bleu,line width=.64pt,line cap=round,line join=round]
    \fill[white] (-.57,-.57) rectangle (.57,.57);
    \draw[rounded corners=1pt,fill=bleu!6] (-.47,-.45) rectangle (.47,.45);
    \draw[rounded corners=1pt] (-.27,-.23) rectangle (.27,.23);
    \foreach \t in {-.31,-.1,.11,.32}{\draw (-.47,\t)--(-.76,\t); \draw (.47,\t)--(.76,\t);}
    \foreach \t in {-.3,-.1,.1,.3}{\draw (\t,-.45)--(\t,-.72);\draw (\t,.45)--(\t,.72);}
    \draw[rouge] (-.27,.14)--(-.1,.14)--(-.1,-.03)--(.14,-.03);
    \fill[rouge] (.14,-.03) circle (.048);
  \end{scope}
  % Antenne, rayonnement radio
  \begin{scope}[shift={([xshift=-8.0mm,yshift=-116mm]current page.north east)},scale=.75,opacity=.27,draw=bleu,line width=.75pt,line cap=round,line join=round]
    \draw (-.43,-.57)--(.45,-.57) (0,-.56)--(0,.5) (-.28,-.34)--(0,.5)--(.28,-.34);
    \fill[rouge] (0,.55) circle (.085);
    \draw[rouge] (.19,.31) arc (42:-42:.35);
    \draw[rouge] (-.19,.31) arc (138:222:.35);
    \draw (.32,.48) arc (42:-42:.62);
    \draw (-.32,.48) arc (138:222:.62);
  \end{scope}
  % Robot a roues, avec capteur et electronique
  \begin{scope}[shift={([xshift=-8.0mm,yshift=-179mm]current page.north east)},scale=.77,opacity=.27,draw=bleu,line width=.75pt,line cap=round,line join=round]
    \draw[rounded corners=2pt,fill=bleu!6] (-.55,-.26) rectangle (.55,.36);
    \draw[fill=rouge!30] (-.35,-.43) circle (.18) (.35,-.43) circle (.18);
    \fill[bleu] (-.35,-.43) circle (.06) (.35,-.43) circle (.06);
    \draw (0,.36)--(0,.57);\fill[rouge] (0,.63) circle (.08);
    \draw (-.24,.13) circle (.09) (.24,.13) circle (.09);
    \draw (-.2,-.08)--(.2,-.08);
    \draw (-.55,.12)--(-.74,.25) (.55,.12)--(.74,.25);
  \end{scope}
  % Travail en equipe : trois silhouettes et un lien vers une tache commune
  \begin{scope}[shift={([xshift=-8.0mm,yshift=-238mm]current page.north east)},scale=.73,opacity=.28,draw=bleu,line width=.75pt,line cap=round,line join=round]
    \draw (0,.32) circle (.14) (-.44,.07) circle (.13) (.44,.07) circle (.13);
    \draw (-.22,-.12) arc (155:25:.28) (-.67,-.48) arc (155:25:.24) (.21,-.48) arc (155:25:.24);
    \draw[rouge] (-.33,-.25)--(0,-.58)--(.33,-.25);
    \draw[fill=rouge!20] (0,-.59) circle (.12);
  \end{scope}
\end{tikzpicture}}
\AddToShipoutPictureBG{\decorationsCV}

% --- Composants typographiques ------------------------------------------------
\titleformat{\section}{\fontsize{9.85}{11}\bfseries\color{bleu}}{}{0pt}{\MakeUppercase}[{\vspace{-2pt}\color{rouge}\titlerule[.66pt]}]
\titlespacing*{\section}{0pt}{6.2pt}{3.1pt}
\setlist[itemize]{leftmargin=11pt,topsep=1.1pt,itemsep=.85pt,parsep=0pt,partopsep=0pt,label={\color{rouge}\raisebox{.4ex}{\tiny$\bullet$}}}
\newcommand{\entry}[4]{%
\begin{tabularx}{\linewidth}{@{}>{\raggedright\arraybackslash}X>{\raggedleft\arraybackslash}p{3.08cm}@{}}
 {\bfseries\color{bleu}#1} & {\fontsize{8.5}{9.3}\selectfont\color{secondaire}#2} \\[-2pt]
 {\fontsize{8.5}{9.6}\selectfont\color{secondaire}#3} & {\fontsize{8.5}{9.3}\selectfont\color{secondaire}#4}
\end{tabularx}\vspace{-3.4pt}\par}
\newcommand{\gapentry}{\vspace{3.1pt}}
\newcommand{\comp}[2]{\textbf{\color{bleu}#1} \hspace{2pt} #2\par}


\begin{document}
% Profil et coordonnees ; elements visuels d'origine conserves
\begin{minipage}[t]{.635\linewidth}
  {\fontsize{21.2}{22.3}\selectfont\bfseries\color{bleu}Youssef EL KATTOUFI}\par\vspace{1.5pt}
  {\fontsize{10.2}{11.5}\selectfont\bfseries\color{rouge}ÉLÈVE-INGÉNIEUR -- ÉLECTRONIQUE ET INFORMATIQUE}\par\vspace{1.1pt}
  {\fontsize{9.15}{10}\selectfont\bfseries\color{bleu}\stage}\par
\end{minipage}\hfill
\begin{minipage}[t]{.355\linewidth}\raggedleft
  {\fontsize{8.7}{10.2}\selectfont
  Choisy-le-Roi, France\\
  \href{tel:+33651654237}{+33 6 51 65 42 37}\\
  \href{mailto:elkattoufi.youssef@hotmail.com}{elkattoufi.youssef@hotmail.com}\\
  \href{https://www.linkedin.com/in/youssef-el-kattoufi-5a04a0282/}{LinkedIn} \quad\textcolor{gris}{|}\quad \href{https://github.com/Darkzveller}{GitHub} }
\end{minipage}
\vspace{4.1pt}
\textcolor{rouge}{\rule{1.35cm}{1.75pt}}\hspace{1pt}\textcolor{gris!65}{\rule{\dimexpr\linewidth-1.45cm\relax}{.6pt}}\par\vspace{3.0pt}
{\fontsize{9.0}{10.7}\selectfont
Élève-ingénieur en Électronique et Informatique à \textbf{Polytech Sorbonne}, je recherche un \textbf{stage dès avril 2027}. Mes expériences m'ont appris à \textbf{développer et tester} des systèmes, à analyser des problèmes concrets et à \textbf{travailler avec plusieurs équipes}. Je souhaite mettre ces acquis en pratique dans les domaines du logiciel, des réseaux ou des systèmes connectés.\par

\section{Expériences professionnelles}
\entry{Apprenti ingénieur -- Génie électrique}{Oct. 2024 -- août 2025}{Syctom -- sites Isséane et Saint-Ouen}{Île-de-France}
\begin{itemize}
  \item \textbf{Sujet :} modernisation d'installations techniques en environnement industriel, avec des échanges entre exploitants, prestataires et équipes de maintenance.
  \item \textbf{Missions et contributions :} analyse des besoins et rédaction de pièces techniques pour le renouvellement de la \textbf{GTB} ; suivi de la modernisation des liaisons de sûreté en \textbf{fibre optique}, commandes et coordination de chantier. Réalisation d'une étude de coût et d'amortissement pour \textbf{6 bornes de recharge}.
\end{itemize}
\gapentry
\entry{Stagiaire technicien électronicien supérieur}{Avr. -- juin 2024}{Cobham Aerospace Communications}{France}
\begin{itemize}
  \item \textbf{Sujet :} fiabiliser et accélérer le contrôle de cartes électroniques. \textbf{Missions :} développement en \textbf{C++} de séquences de tests automatisés, création d'une interface de commande, réalisation d'essais et diagnostic des anomalies.
  \item \textbf{Livrables et outils :} banc de test, documentation technique et simulations sous \textbf{LTspice}. \textbf{Objectif du projet :} réduire la durée de test de \textbf{30\,\%} (objectif, non résultat mesuré).
\end{itemize}
\gapentry
\entry{Surveillant sauveteur aquatique}{Juil. -- août 2026}{Centre aquatique de Villecresnes -- GPSEA}{Villecresnes}
\begin{itemize}
  \item Surveillance et prévention des risques dans un établissement accueillant jusqu'à \textbf{1\,000 personnes} : coordination avec les autres sauveteurs, communication avec le public et interventions de premiers secours.
\end{itemize}

\section{Projets techniques \& associatifs}
\entry{Coupe de France de robotique -- Responsable électronique / asservissement}{2024 -- 2025}{Club CRAC -- projet collectif de robot autonome}{Associatif}
\begin{itemize}
  \item \textbf{Sujet :} faire fonctionner ensemble les cartes et la motorisation du robot. \textbf{Missions :} développement du socle \textbf{C++} sur \textbf{ESP32/STM32}, commandes des actionneurs, échanges \textbf{CAN}, interface de tests et mises à jour \textbf{OTA}.
  \item \textbf{Contribution :} commandes de \textbf{2 moteurs} avec odométrie et asservissements \textbf{PID} en vitesse/position ; intégration électronique et logicielle avec l'équipe. \textbf{Outils :} FreeRTOS, Git, PlatformIO.
\end{itemize}
\gapentry
\entry{Projet IoT -- Acquisition et échange de mesures}{2026}{Polytech Sorbonne -- capteurs connectés}{Universitaire}
\begin{itemize}
  \item \textbf{Sujet / réalisation :} récupérer des mesures et les rendre accessibles à d'autres applications. Acquisition sur \textbf{ESP32}, transmission par \textbf{MQTT} et \textbf{HTTP REST}, suivi des données sous \textbf{Node-RED}.
\end{itemize}
\gapentry
\entry{Trottinette à double motorisation -- moteur sans terres rares}{2024 -- 2025}{IUT de Cachan -- conception et électronique de puissance}{Universitaire}
\begin{itemize}
  \item \textbf{Sujet / missions :} concevoir le pilotage d'un moteur \textbf{AC} et d'un moteur \textbf{DC} d'assistance. Conception d'une \textbf{carte PCB}, commandes \textbf{STM32/ESP32}, suivi des températures, des phases et de la consommation ; \textbf{CAO 3D} et essais d'onduleur triphasé.
\end{itemize}
\gapentry
\entry{Traitement du signal -- Réception et décodage DCF77}{2025 -- 2026}{Polytech Sorbonne -- signal radio bruité à 77,5 kHz}{Universitaire}
\begin{itemize}
  \item \textbf{Problème / solution :} retrouver une information binaire dans un signal radio bruité. Analyse à l'oscilloscope et WebSDR, filtre RII d'ordre 8, traitement \textbf{MATLAB} avec \textbf{filtfilt}, extraction de l'enveloppe et seuillage adaptatif ; identification des impulsions de \textbf{100 et 200 ms}.
\end{itemize}

\section{Formation}
\entry{Diplôme d'ingénieur -- Électronique et Informatique}{2025 -- 2028}{Polytech Sorbonne -- Sorbonne Université}{Paris}
\vspace{1pt}
\entry{BUT Génie électrique et informatique industrielle -- option Électricité et maîtrise de l'énergie}{2022 -- 2025}{IUT de Cachan -- Université Paris-Saclay}{Cachan}

\section{Compétences mises en pratique}
{\fontsize{8.65}{10.05}\selectfont
\comp{Résolution de problèmes :}{diagnostic et essais chez Cobham ; traitement du signal sous MATLAB (DCF77).}
\comp{Travail en équipe / organisation :}{coordination d'intervenants au Syctom ; intégration du robot avec le club CRAC ; prévention en équipe au centre aquatique.}
\comp{Outils complémentaires :}{Python, KiCad, SolidWorks, STM32CubeIDE, UART, SPI, I\textsuperscript{2}C, Git ; BNSSA, PSE2 ; français, anglais B1.}
}
\end{document}
